% ---------------------------------------------------------------------------
% Networked communication environments: LineMsg and WirelessComm
% ---------------------------------------------------------------------------
\begingroup
% Line breaking for paragraphs with long, unbreakable identifiers (local to
% this chapter): allow looser lines instead of overfull ones.
\setlength{\emergencystretch}{3em}\tolerance=1000
\chapter{Networked Communication Environments}\label{chap:networked}

This chapter documents two discrete multi-agent benchmarks in which agents on a
network cooperate under local observations:

\begin{itemize}
  \item \textbf{LineMsg} (\code{env\_lib.linemsg\_env.LineMsgEnv}): agents on
        a line relay a message from a source to a sink;
  \item \textbf{WirelessComm}
        (\code{env\_lib.wireless\_comm\_env.WirelessCommEnv}): agents on a grid
        send deadline-constrained packets to shared access points and must avoid
        collisions.
\end{itemize}

Both environments have small discrete state and action spaces, local
observation windows controlled by \code{n\_obs\_neighbors}, per-agent rewards
and a fixed episode length. They are designed for research on decentralised
and networked multi-agent reinforcement learning: exploiting the locality of
interactions (policies and value functions that depend on a $\kappa$-hop
neighbourhood), credit assignment, communication and coordination protocols,
scalability in the number of agents, and robustness to failing agents. The
WirelessComm environment is the Gymnasium joint-action port of an earlier
PettingZoo implementation used in the DSDP project.

% ===========================================================================
\section{LineMsg}\label{sec:linemsg}

% ---------------------------------------------------------------------------
\subsection{Model}\label{sec:linemsg-model}

$N$ agents sit on a line. Agent $N-1$ at the right end is the source, agent 0
at the left end is the sink. Each agent holds a binary state $s_i \in \{0, 1\}$
(1: it holds the message) and chooses a binary action $a_i \in \{0, 1\}$ at
every step. All states are updated synchronously from the previous states:
\begin{align}
  s_0 &\leftarrow s_1, \nonumber\\
  s_i &\leftarrow a_i \wedge \bigl(s_{i+1} \vee [u_i < 0.8]\bigr),
       \qquad 0 < i < N-1,\; u_i \sim \mathcal{U}(0, 1), \label{eq:linemsg}\\
  s_{N-1} &\leftarrow a_{N-1}. \nonumber
\end{align}
An active middle agent ($a_i = 1$) therefore receives the message from its
right neighbour for sure and otherwise obtains one with probability 0.8; an
inactive agent loses its message. The sink copies agent 1 and its own action
has no effect; the source holds the message exactly when it acts. At reset all
agents hold the message ($s_i = 1$).

\paragraph{Reward.} Agent $i$ receives
\begin{equation}
  r_i = w_i\, s_i, \qquad w_0 = 1.0, \quad w_i = 0.1 \;\; (i > 0),
\end{equation}
evaluated on the new states, and the team reward is $R = \sum_i r_i$. The
maximum per step is $1 + 0.1(N - 1)$, which the ``always relay'' policy
($a_i = 1$ for all $i$) attains at every step.

\paragraph{Episode end.} \code{terminated} is always \code{False};
\code{truncated} becomes \code{True} after \code{max\_iter} steps.

% ---------------------------------------------------------------------------
\subsection{Spaces}\label{sec:linemsg-spaces}

\paragraph{Action.} The joint action can be given in two formats, both of
which \code{step()} always accepts:
\begin{itemize}
  \item an integer in \code{Discrete(2**N)} whose bit $i$ is the action of
        agent $i$ (the declared space for \code{action\_space\_type="discrete"},
        the default; only possible for $N \le 62$);
  \item an array of $N$ values in $\{0, 1\}$ (integer or boolean), the
        declared space \code{MultiBinary(N)} for
        \code{action\_space\_type="multibinary"}, which has no limit on $N$.
\end{itemize}
Out-of-range integers, arrays of the wrong shape and values other than 0 and 1
raise \code{ValueError}.

\paragraph{Observation.} \code{Box(0, 3, (N, 2k+1), float32)} with
$k$ = \code{n\_obs\_neighbors}. The state vector is padded with $k$ boundary
cells on each side, and row $i$ is the window of the $2k + 1$ cells centred on
agent $i$:

\begin{table}[htbp]
  \centering\small
  \caption{LineMsg observation (row $i$).}
  \label{tab:linemsg-obs}
  \begin{tabularx}{\textwidth}{@{}l>{\raggedright\arraybackslash}X@{}}
    \toprule
    Column & Content \\
    \midrule
    \code{0 : k} & states of agents $i-k, \dots, i-1$ (left neighbours, towards the sink) \\
    \code{k} & own state $s_i$ \\
    \code{k+1 : 2k+1} & states of agents $i+1, \dots, i+k$ (right neighbours, towards the source) \\
    \midrule
    Values & 0: no message, 1: message, 2: boundary cell beyond the end of the line \\
    \bottomrule
  \end{tabularx}
\end{table}

% ---------------------------------------------------------------------------
\subsection{Parameters and info}\label{sec:linemsg-params}

\begin{table}[htbp]
  \centering\small
  \caption{Constructor parameters of \code{LineMsgEnv}.}
  \label{tab:linemsg-params}
  \begin{tabularx}{\textwidth}{@{}>{\raggedright\arraybackslash}p{0.25\textwidth}>{\raggedright\arraybackslash}p{0.14\textwidth}>{\raggedright\arraybackslash}X@{}}
    \toprule
    Parameter & Default & Description \\
    \midrule
    \code{num\_agents} & \code{10} & Number of agents $N$ ($\ge 2$: a sink and a source, plus any number of relays). \\
    \code{n\_obs\_neighbors} & \code{1} & Observation radius $k$; $0$ is treated as $1$ with a warning. \\
    \code{max\_iter} & \code{50} & Episode length; reaching it sets \code{truncated}. \\
    \code{render\_mode} & \code{None} & \code{None}, \code{"human"}, \code{"rgb\_array"} or \code{"ansi"}. \\
    \code{action\_space\_type} & \code{"discrete"} & Declared action space: \code{"discrete"} (\code{Discrete(2**N)}, $N \le 62$) or \code{"multibinary"} (\code{MultiBinary(N)}). \\
    \bottomrule
  \end{tabularx}
\end{table}

The info dictionary of \code{reset()} and \code{step()} contains a single key,
\code{"agent\_rewards"}: the \code{float64} per-agent rewards $r_i$ (zeros at
reset). The attribute \code{state} holds the padded state vector (agent $i$ is
\code{state[k + i]}), \code{actions} the last joint action with the same padding
(2 before the first step), and \code{num\_moves} the number of steps taken.
\code{step()} before \code{reset()} raises \code{env\_lib.ResetNeededError}.

% ---------------------------------------------------------------------------
\subsection{Rendering}\label{sec:linemsg-render}

The dashboard (Figure~\ref{fig:linemsg}) shows the line of agents at the top,
coloured by state, with a halo under agents that act and arrows on the active
links, which point in the direction of message flow towards the sink. Below it
a space-time raster (one row per step, one column per agent) makes message
fronts visible, and the right panel traces the team reward with its maximum.
The header reports the step, the last reward, the episode return and the number
of agents holding the message. \code{render\_mode="ansi"} returns a text line
with the action and state of every agent.

\begin{figure}[htbp]
  \centering
  \includegraphics[width=\figwidth]{linemsg.png}
  \caption{LineMsg dashboard (12 agents, step 30 of 50).
    Top: agent states and active links. Bottom left: space-time raster of the
    agent states. Bottom right: team reward per step.}
  \label{fig:linemsg}
\end{figure}

% ---------------------------------------------------------------------------
\subsection{Usage}\label{sec:linemsg-usage}

\begin{pycode}
import numpy as np
from env_lib.linemsg_env import LineMsgEnv

env = LineMsgEnv(num_agents=10, action_space_type="multibinary")
obs, info = env.reset(seed=0)                 # obs: (10, 3) float32
rng = np.random.default_rng(0)
episode_return = 0.0
while True:
    action = (rng.random(env.num_agents) < 0.9).astype(np.int64)  # duty cycle
    obs, reward, terminated, truncated, info = env.step(action)
    episode_return += reward
    if terminated or truncated:
        break
print(episode_return, info["agent_rewards"])

env.reset(seed=1)
env.step(0b1111111110)          # integer action: all agents but the sink act
\end{pycode}

The example script compares a random policy, duty-cycled relays and the
always-relay policy and can record an episode:

\begin{shellcode}
python examples/environments/linemsg_demo.py --policy relay --save renders/linemsg.gif
python examples/environments/linemsg_demo.py --num-agents 30 --episodes 50
python examples/environments/linemsg_demo.py --render-mode ansi
\end{shellcode}

Over 20 episodes with $N = 10$ and \code{max\_iter=50} the mean returns are
44.1 for uniformly random actions, 85.2 for duty-cycled relays that act with
probability 0.9, and 95.0 for the always-relay policy (the maximum).

% ===========================================================================
\section{WirelessComm}\label{sec:wireless}

% ---------------------------------------------------------------------------
\subsection{Model}\label{sec:wireless-model}

$G_x \times G_y$ agents sit on a grid (\code{grid\_x} rows, \code{grid\_y}
columns); agent $(i, j)$ has index $i G_y + j$. Access points sit at the
$(G_x - 1)(G_y - 1)$ interior corners of the grid: access point $(a, b)$ lies
between agents $(a, b)$, $(a+1, b)$, $(a, b+1)$ and $(a+1, b+1)$. Every agent
holds a packet queue with $D$ = \code{ddl} deadline slots; slot 0 holds the
packet with the earliest deadline. Each agent chooses one of five actions per
step (Table~\ref{tab:wireless-actions}).

\begin{table}[htbp]
  \centering\small
  \caption{WirelessComm actions of agent $(i, j)$. A direction without an
    access point (at the border of the grid) is a failed attempt.}
  \label{tab:wireless-actions}
  \begin{tabular}{@{}cll@{}}
    \toprule
    Action & Meaning & Access point \\
    \midrule
    0 & stay idle & none \\
    1 & transmit up-left & $(i-1, j-1)$ \\
    2 & transmit down-left & $(i, j-1)$ \\
    3 & transmit up-right & $(i-1, j)$ \\
    4 & transmit down-right & $(i, j)$ \\
    \bottomrule
  \end{tabular}
\end{table}

\paragraph{Transition.} Let $\ell_m$ be the load of access point $m$, the
number of agents that transmit to it in the step; every transmitting agent
counts, whether or not it holds a packet. The transmission of agent $k$ is
\emph{eligible} when the agent transmits, holds at least one packet, its access
point exists and that access point has load exactly $\ell_m = 1$. Each eligible
transmission succeeds with probability $q$; one uniform draw is made per
eligible agent, in agent order. A success removes the agent's earliest-deadline
packet and earns reward 1. Afterwards every queue shifts one slot towards the
deadline (packets in slot 0 expire and are lost), and a new packet arrives in
the last slot with probability $p$. At reset every queue slot is occupied
independently with probability $1/2$.

\paragraph{Reward and outcomes.} The per-agent reward is 1 for a delivered
packet and 0 otherwise; the team reward is the number of packets delivered in
the step. The outcome of every agent is reported in \code{info["outcomes"]}
(Table~\ref{tab:wireless-info}): idle (0), success (1), collision (2: the
access point exists and has load $\ge 2$) or lost (3: any other failed attempt,
that is no packet, no access point or a channel loss).

\paragraph{Episode end.} \code{terminated} is always \code{False};
\code{truncated} becomes \code{True} after \code{max\_iter} steps.

% ---------------------------------------------------------------------------
\subsection{Spaces}\label{sec:wireless-spaces}

\paragraph{Action.} \code{MultiDiscrete([5] * n)} with $n = G_x G_y$. An action
of shape \code{(n,)} or \code{(grid\_x, grid\_y)} is accepted; other shapes and
values outside $0, \dots, 4$ raise \code{ValueError}.

\paragraph{Observation.} \code{Box(0, 2, (n, D w\^{}2), float32)} with
$w = 2k + 1$ and $k$ = \code{n\_obs\_neighbors}. The queues form a
$(D, G_x, G_y)$ array that is padded with $k$ cells of value 2 on every side of
the grid. Row $m$ of the observation is the $(D, w, w)$ window centred on agent
$m$, flattened in C order (slot, then row, then column). Values are 0 (empty
slot), 1 (packet) and 2 (outside the grid). The agent's own slot $d$ is column
$d w^2 + k w + k$. The default configuration ($D = 2$, $k = 1$) has 18 columns.

% ---------------------------------------------------------------------------
\subsection{Parameters and info}\label{sec:wireless-params}

\begin{table}[htbp]
  \centering\small
  \caption{Constructor parameters of \code{WirelessCommEnv}.}
  \label{tab:wireless-params}
  \begin{tabularx}{\textwidth}{@{}>{\raggedright\arraybackslash}p{0.39\textwidth}>{\raggedright\arraybackslash}p{0.1\textwidth}>{\raggedright\arraybackslash}X@{}}
    \toprule
    Parameter & Default & Description \\
    \midrule
    \code{grid\_x}, \code{grid\_y} & \code{6}, \code{6} & Number of agent rows and columns (each $\ge 2$, so that at least one access point exists). \\
    \code{ddl} & \code{2} & Deadline horizon $D$: queue slots per agent ($\ge 1$). \\
    \code{packet\_arrival\_probability} & \code{0.8} & Arrival probability $p$ of a new packet per agent and step. \\
    \code{success\_transmission\_probability} & \code{0.8} & Success probability $q$ of an eligible transmission. \\
    \code{n\_obs\_neighbors} & \code{1} & Observation radius $k$ (Chebyshev distance, $\ge 0$). \\
    \code{max\_iter} & \code{50} & Episode length; reaching it sets \code{truncated}. \\
    \code{render\_mode} & \code{None} & \code{None}, \code{"human"}, \code{"rgb\_array"} or \code{"ansi"}. \\
    \bottomrule
  \end{tabularx}
\end{table}

\begin{table}[htbp]
  \centering\small
  \caption{WirelessComm info dictionary (the same keys, all zero, at reset).}
  \label{tab:wireless-info}
  \begin{tabularx}{\textwidth}{@{}l>{\raggedright\arraybackslash}X@{}}
    \toprule
    Key & Content \\
    \midrule
    \code{agent\_rewards} & \code{float64} array \code{(n,)}: 1 for a delivered packet, else 0 \\
    \code{outcomes} & \code{int8} array \code{(n,)}: 0 idle, 1 success, 2 collision, 3 lost
      (constants \code{OUTCOME\_IDLE}, \code{OUTCOME\_SUCCESS}, \code{OUTCOME\_COLLISION},
      \code{OUTCOME\_LOST} and \code{OUTCOME\_LABELS} in \code{env\_lib.wireless\_comm\_env.wireless\_comm\_env}) \\
    \code{ap\_load} & \code{int64} array \code{(grid\_x - 1, grid\_y - 1)}: number of agents that transmitted to each access point \\
    \bottomrule
  \end{tabularx}
\end{table}

The environment also exposes \code{last\_outcomes} and \code{last\_ap\_load}
(the values of the last step), \code{access\_point\_mapping(i, j, action)},
which returns the access point $(a, b)$ reached by agent $(i, j)$ or
\code{(None, None)}, and \code{check\_transmission\_fail(a, b, load)}.
\code{step()} before \code{reset()} raises \code{env\_lib.ResetNeededError}.

% ---------------------------------------------------------------------------
\subsection{Rendering}\label{sec:wireless-render}

The dashboard (Figure~\ref{fig:wireless}) draws the agent grid on the left.
Every agent is a circle coloured by the outcome of its last action, with a
small queue bar below it (one cell per deadline slot; the left cell expires
first). Access points are squares at the interior corners, coloured by their
load (idle, one transmitter, collision), and a line joins every transmitting
agent to the access point it chose, coloured by the outcome; the line of an attempt
towards a missing access point ends outside the grid. The right panels show the packets
delivered per step with a rolling mean and the rolling collision count, and the
cumulative return. \code{render\_mode="ansi"} returns the action and queue of
every agent as text.

\begin{figure}[htbp]
  \centering
  \includegraphics[width=\figwidth]{wireless.png}
  \caption{WirelessComm dashboard ($6 \times 6$ agents, 25 access points,
    step 25 of 50). Left: agents coloured by outcome with their packet queues,
    access points coloured by load and transmission links. Right: packets
    delivered and collisions per step (10-step means) and the cumulative
    return.}
  \label{fig:wireless}
\end{figure}

% ---------------------------------------------------------------------------
\subsection{Usage}\label{sec:wireless-usage}

The following policy is a collision-free rotating schedule: at step $t$ every
agent that holds a packet uses the same direction $(4, 1, 2, 3)[t \bmod 4]$ if
an access point exists in that direction. A common direction maps agents to
access points one-to-one, so no collisions occur.

\begin{pycode}
import numpy as np
from env_lib.wireless_comm_env import WirelessCommEnv

env = WirelessCommEnv(grid_x=6, grid_y=6)
obs, info = env.reset(seed=0)                     # obs: (36, 18) float32
k, w = env.n_obs_nghbr, env.window
own = [d * w * w + k * w + k for d in range(env.ddl)]  # own queue columns
directions = (4, 1, 2, 3)
cells = [(i, j) for i in range(env.grid_x) for j in range(env.grid_y)]
reachable = {d: np.array([env.access_point_mapping(i, j, d)[0] is not None
                          for i, j in cells])
             for d in directions}
episode_return = 0.0
for t in range(env.max_iter):
    d = directions[t % 4]
    has_packet = obs[:, own].max(axis=1) > 0.5
    action = np.where(has_packet & reachable[d], d, 0)
    obs, reward, terminated, truncated, info = env.step(action)
    episode_return += reward
print(episode_return, np.bincount(info["outcomes"], minlength=4))
\end{pycode}

The example script compares uniformly random actions, $p$-persistent slotted
ALOHA and the rotating schedule:

\begin{shellcode}
python examples/environments/wireless_comm_demo.py --policy schedule \
    --save renders/wireless.gif
python examples/environments/wireless_comm_demo.py --grid 8 8 --tau 0.5 --episodes 50
\end{shellcode}

Over 20 episodes on the $6 \times 6$ grid the mean returns are 385 for random
actions, 398 for ALOHA with transmit probability 0.6 and 910 for the rotating
schedule, which causes no collisions.

% ===========================================================================
\section{Registered ids}\label{sec:networked-ids}

\begin{table}[htbp]
  \centering\small
  \caption{Registered ids of the networked environments.}
  \label{tab:networked-ids}
  \begin{tabular}{@{}lll@{}}
    \toprule
    Id & Class & Configuration \\
    \midrule
    \envid{LineMsg-v0} & \code{LineMsgEnv} & defaults (10 agents) \\
    \envid{WirelessComm-v0} & \code{WirelessCommEnv} & defaults ($6 \times 6$ grid, 36 agents) \\
    \envid{WirelessComm-v1} & \code{WirelessCommEnv} & \code{grid\_x=4}, \code{grid\_y=4} (16 agents) \\
    \bottomrule
  \end{tabular}
\end{table}

% ===========================================================================
\section{Reproducibility and performance}\label{sec:networked-perf}

Both environments draw all randomness (relay successes in LineMsg; initial
queues, transmission successes and packet arrivals in WirelessComm) from
\code{self.np\_random}, seeded by \code{reset(seed=...)}. Seeded trajectories
are bit-identical to those of the implementation before version 1.0, so
earlier results can be reproduced exactly.

Measured on the development machine without rendering, a LineMsg step with 10
agents takes 0.013\,ms (0.019\,ms before 1.0), and a WirelessComm step takes
0.05\,ms on the $6 \times 6$ grid (0.154\,ms before) and 0.046\,ms on a
$12 \times 12$ grid (0.56\,ms before). Rendering a dashboard frame in
\code{"rgb\_array"} mode takes about 15\,ms for either environment; the history
needed by the dashboards is only recorded in the \code{"human"} and
\code{"rgb\_array"} modes.

% ===========================================================================
\section{Changes in 1.0}\label{sec:networked-changes}

The seeded dynamics of both environments are unchanged. Interface changes are
summarised below; see Appendix~\ref{app:migration} for migration advice.

\paragraph{LineMsg.}
\begin{itemize}
  \item Action decoding and observations are vectorised.
  \item The joint action may be an integer or an array of $N$ bits. The new
        option \code{action\_space\_type="multibinary"} declares a
        \code{MultiBinary(N)} space and supports more than 62 agents.
  \item \code{step()} before \code{reset()} raises instead of resetting
        implicitly.
  \item The former text output is \code{render\_mode="ansi"}; \code{"human"}
        and \code{"rgb\_array"} show the new dashboard.
\end{itemize}

\paragraph{WirelessComm.}
\begin{itemize}
  \item \code{step()} is fully vectorised (up to 37 times faster on large grids).
  \item Fixed: the stored actions were never updated, so rendering showed no
        actions. Per-agent outcomes and access-point loads are reported in
        \code{info}.
  \item Arguments and actions are validated; \code{"ansi"} text mode and a new
        dashboard were added.
  \item \code{env\_lib.wireless\_comm\_env.register()} is deprecated; the ids
        are registered when \pkg{env\_lib} is imported.
\end{itemize}

\endgroup
